\section{The actual algebraic weight representation with all integral fields}
\label{sec:weight-representation-v42}

\status{This chapter constructs the existential weight representation
used by the original \texttt{ProdZpowConst.wtRep}, on the specified
algebraic representation and with the specified field isomorphism.
It proves all ten assertion fields of \texttt{WtRep}.  The proof uses
the already reconstructed characteristic-zero polynomial results
HP1--HP2, and a saturated lattice over the actual coefficient DVR.
It therefore does not need a new general standard-monomial basis
theorem over arbitrary domains.  The source endpoint permits the
basis chosen here: its required identification is an explicit
intertwining isomorphism with the unchanged original
\texttt{algRep}.}

Fix a number field $L$, a prime $\ell$, an integer $n\geq0$, a
finite index set $V$, places $S(v)$ of $L$, a field isomorphism
\[
             \iota:\overline{\mathbb Q}_\ell\xrightarrow{\sim}
                                                              \mathbb C,
\]
and the original weight function $\lambda$ assigning a dominant
integer weight of rank $n$ to each embedding
$L\hookrightarrow\overline{\mathbb Q}_\ell$.  The isomorphism
$\iota$ is algebraic and is never assumed continuous.  Let $c$
be the original coefficient datum: its ring $c.O$ is a DVR
subalgebra of $\overline{\mathbb Q}_\ell$, finite over
$\mathbb Z_\ell$.

Put $F_v=L_{S(v)}$ and let
\[
 \mathcal E=\{(v,j):v\in V,\ j:F_v\hookrightarrow
           \overline{\mathbb Q}_\ell\text{ a continuous field map}\}.
\]
This is the literal source \texttt{Emb}.  For $e=(v,j)$ write
$\lambda_e=\lambda(j|_L)$ and $\psi_e=\iota\circ j$.
The original group and representation are
\begin{equation}\label{eq:wrp-algrep}
 G=\prod_{v\in V}\operatorname{GL}_n(F_v),\qquad
 \rho(g)=\bigotimes_{e=(v,j)\in\mathcal E}
                 \operatorname{sqRep}(\psi_e(g_v),\lambda_e)
\end{equation}
on the actual tensor product
\[
 \mathcal V_{\mathbb C}=
       \bigotimes_{e\in\mathcal E}
                  \operatorname{nablaSq}_{\mathbb C}(n,\lambda_e).
\]
All indices $i$ below run from $0$ to $n-1$, as in the source.
For a dominant weight $\lambda_e$, set $s_e=(\lambda_e)_{n-1}$
and $\nu_{e,i}=(\lambda_e)_i-s_e$ if $n>0$; for $n=0$ use
$s_e=0$ and the constant polynomial.  By HP1, its literal
polynomial factor is
\begin{equation}\label{eq:wrp-polynomial}
 \mathcal N_{\nu_e}(K)=\{f:f(Xt)=t^{\nu_e}f(X),\
                                  f(Xu)=f(X)\ (u\in N_+)\},
 \qquad \rho_e(g)f=\det(g)^{s_e}f(g^{\mathsf T}X).
\end{equation}
The covariance conditions are polynomial identities, or the
finite integer-coefficient kernel equations of HP1.  The
determinant power may be negative.

The local Iwahori datum is the original $\mathrm{iwS}(n,S)$.
We use its actual uniformizers $\varpi_v$, not new uniformizers
chosen after extending coefficients.  OST26--OST27 identify its
torus with diagonal matrices, its compact groups with integral
matrices, and its positive monoid with tuples
\[
      a_v=\operatorname{diag}(\varpi_v^{a_{v,0}},\ldots,
                             \varpi_v^{a_{v,n-1}}),
                    \quad a_{v,0}\geq\cdots\geq a_{v,n-1}.
\]
Its distinguished $z$ uses $a_{v,i}=n-1-i$ at every $v$.

\subsection{One coefficient extension for every local embedding}

% BEGIN WRP1
\paragraph{WRP1. The original embedding index is finite and its images are finite local fields.}
\label{claim:WRP1}
The set $\mathcal E$ is finite.  Every $j$ occurring in it has
$S(v)$ above $\ell$, fixes the embedded $\mathbb Q_\ell$, and
has image a finite extension of $\mathbb Q_\ell$ inside the
specified $\overline{\mathbb Q}_\ell$.  Moreover
\begin{equation}\label{eq:wrp-local-integrality}
 |j(x)|_\ell\leq1\quad(x\in\mathcal O_{F_v}),\qquad
 |j(x)|_\ell<1\quad(x\in\mathfrak m_{F_v}).
\end{equation}
\begin{proof}
Continuous maps agreeing on the dense subfield $L$ agree on its
completion.  There are only finitely many field maps
$L\to\overline{\mathbb Q}_\ell$: for a finite set of algebraic
generators each image is a root of a fixed nonzero polynomial,
so there are finitely many possible lists of images.  This and
finiteness of $V$ prove the first assertion.

If $S(v)$ has residue characteristic $p\ne\ell$, then
$p^r\to0$ in $F_v$, whereas its image under $j$ has
$\ell$-adic norm $1$ for every $r$, contradicting continuity.
Thus $S(v)$ lies above $\ell$.  The inclusion of $\mathbb Q$
extends to the usual inclusion $\mathbb Q_\ell\subset F_v$;
its topology is obtained by restricting the local valuation,
and completion supplies this inclusion.  By continuity and
density of $\mathbb Q$ in $\mathbb Q_\ell$, $j$ is the identity
on that subfield.  The completion $F_v$ is finite over
$\mathbb Q_\ell$: the $\mathbb Q_\ell$-span of a rational
basis of $L$ contains the dense subfield $L$ and is closed.
To check closedness, use the coordinate compactness argument
of GA2: a linear injection from a finite dimensional
$\mathbb Q_\ell$-space into a normed space has equivalent
coordinate and induced norms, since the coordinate unit
sphere is compact and its image avoids zero.  Hence its
image is complete and closed.  Consequently this span is
the whole $F_v$.  The $\mathbb Q_\ell$-linear injection $j$
therefore has finite dimensional image, a field.

The local integer ring $\mathcal O_{F_v}$ is compact.
Continuity bounds $|j(x)|_\ell$ on it by some constant.
For any fixed $x$ in that ring all positive powers remain
in the ring, so $|j(x)|_\ell^r$ stays bounded.  This forces
$|j(x)|_\ell\leq1$.  If $x\in\mathfrak m_{F_v}$ then
$x^r\to0$, and continuity forces $j(x)^r\to0$, hence
$|j(x)|_\ell<1$.  These arguments use no continuity of
$\iota$.
\end{proof}
% END WRP1

% BEGIN WRP2
\paragraph{WRP2. One actual coefficient DVR contains all local integer images.}
\label{claim:WRP2}
There is an original coefficient datum $c'$ with
$c.O\subset c'.O\subset\overline{\mathbb Q}_\ell$ such that
every $j$ in WRP1 restricts to a local homomorphism
\[
               \phi_e:\mathcal O_{F_v}\longrightarrow c'.O.
\]
In particular $\phi_e(\mathfrak m_{F_v}^k)
\subset\mathfrak m_{c'.O}^k$ for every $k$.
\begin{proof}
Let $E$ be the compositum inside the specified algebraic closure
of $\operatorname{Frac}(c.O)$ and the finitely many fields
$j(F_v)$.  It is finite over $\mathbb Q_\ell$.  For the first
field this follows from finite $\mathbb Z_\ell$-generation of
$c.O$: as proved in GA2, every denominator in that DVR divides
a power of $\ell$, so its fraction field is spanned by the
same generators over $\mathbb Q_\ell$.  The other fields are
finite by WRP1, and a finite compositum of finite extensions
is finite, by adjoining their finite bases successively.

Define the literal subring
\[
                       O'=\{x\in E:|x|_\ell\leq1\}.
\]
Here is a direct verification that it is an allowed coefficient
DVR.  Choose a $\mathbb Q_\ell$-basis of $E$.  The coordinate
norm and the induced norm are equivalent by the compact unit
sphere argument in WRP1.  Consequently $O'$ is contained
in $\ell^{-M}\mathbb Z_\ell^r$ and contains
$\ell^M\mathbb Z_\ell^r$ for some $M$.  It is a
$\mathbb Z_\ell$-submodule, hence finite over the Noetherian
PID $\mathbb Z_\ell$, and it spans $E$.  The finite algebra
$O'$ is Noetherian: every ideal is a submodule of its finite
$\mathbb Z_\ell$-module, and a finite list of generators
over $\mathbb Z_\ell$ also generates it as an ideal.

Its unique maximal ideal is
$\mathfrak m' =\{x:|x|_\ell<1\}$: an element of norm one
has its inverse in $O'$, while products involving an element
of norm less than one cannot equal one.  This ideal is
nonzero, since it contains $\ell$.  Choose a finite set of
its generators and let $\pi'$ be a nonzero generator of
largest norm.  Every $x\in\mathfrak m'$ is a sum of these
generators with coefficients of norm at most one, so
$|x|_\ell\leq|\pi'|_\ell<1$.  Thus $x/\pi'\in O'$,
and $\mathfrak m'=(\pi')$.  For any nonzero $x\in O'$,
choose the largest $r\geq0$ with
$|x|_\ell\leq|\pi'|_\ell^r$.  Such an $r$ exists and is
finite.  The element $x/(\pi')^r$ is in $O'$ and is a unit;
otherwise its norm would be at most $|\pi'|_\ell$, contrary
to maximality of $r$.  Hence every nonzero element is a unit
times a power of $\pi'$, and taking the smallest exponent
in any nonzero ideal proves the DVR property.  The ring is
complete: $E$ is complete by norm equivalence, $O'$ is
closed, and the powers of $\pi'$ give its norm neighborhoods.
Its residue field is finite because it is a quotient of the
finite ring $O'/\ell O'$.

Every element of $c.O$ has norm at most one.  Indeed finite
generation over $\mathbb Z_\ell$ gives a monic equation by
the determinant trick, and a root with norm greater than one
would give a uniquely largest term in that equation.  Thus
$c.O\subset O'$.  WRP1 gives $j(\mathcal O_{F_v})\subset O'$
and $j(\mathfrak m_{F_v})\subset\mathfrak m'$.  Set
$c'.O=O'$ as a subalgebra of the same algebraic closure.
The asserted maps are the literal restrictions of $j$;
they are local, and the assertion for powers follows by
expanding finite sums of products.  This constructs the
whole coefficient datum, not just an abstract isomorphic
DVR.  Its complex algebra map is always $\iota|_{O'}$.
\end{proof}
% END WRP2

\subsection{Integral weight coordinates on the unchanged polynomial factors}

% BEGIN WRP3
\paragraph{WRP3. A saturated integral lattice stable under all integral matrices.}
\label{claim:WRP3}
Write $O=c'.O$, $K=\operatorname{Frac}(O)$ and fix one dominant
weight $\lambda_e$.  Let $P_O$ be the finite free module of
polynomials of its prescribed column multidegree, and put
\[
              L_e=\mathcal N_{\nu_e}(K)\cap P_O.
\]
Then $L_e$ is a finite free $O$-lattice spanning
$\mathcal N_{\nu_e}(K)$.  It is stable under the original
$\rho_e(g)$ for $g\in\operatorname{GL}_n(O)$.  If
$g\equiv1\pmod{\mathfrak m_O^k}$, then
\begin{equation}\label{eq:wrp-congruence}
             (\rho_e(g)-1)L_e\subset\mathfrak m_O^k L_e.
\end{equation}
\begin{proof}
The submodule $L_e\subset P_O$ is finite and torsion free
over a DVR, hence free.  For clarity this module fact follows
by induction on the rank of $P_O$: project to its first
coordinate, whose nonzero image is a principal ideal,
choose a lift of its generator, and split off that free
summand; apply induction to the kernel in the remaining
coordinates.  A basis of the $K$-space has finitely many
polynomial coefficients; clearing their denominators places
scalar multiples in $L_e$, so it spans the whole space.
By the finite integer-coefficient equations in HP1, this
intersection is also the literal integral polynomial kernel:
an equation over $O$ holds if and only if it holds after
the injective map $O\to K$.

For $g\in\operatorname{GL}_n(O)$, matrix substitution
preserves integral polynomial coefficients and each column
degree.  The multiplier $\det(g)^{s_e}$ is an element of
$O^\times$, including for negative $s_e$.  Left substitution
commutes with the right covariance conditions defining
$\mathcal N_{\nu_e}$, so it preserves that subspace as well.
The same argument for $g^{-1}$ gives stability of the lattice
in both directions.

If $g\equiv1$ modulo $\mathfrak m_O^k$, reduction of matrix
substitution is the identity.  Also
$\det(g)^{s_e}\equiv1$ modulo that ideal; negative powers
cause no difficulty since the determinant is a unit and
$u^{-1}-1=-u^{-1}(u-1)$.  Thus the difference in
\eqref{eq:wrp-congruence} lies in
$\mathfrak m_O^kP_O\cap\mathcal N_{\nu_e}(K)$.
Writing $\mathfrak m_O^k=(\pi'^k)$ and dividing the
polynomial by $\pi'^k$ identifies this intersection with
$\pi'^kL_e$.  This proves the exact congruence, without
assuming an integral standard-monomial theorem.
\end{proof}
% END WRP3

% BEGIN WRP4
\paragraph{WRP4. A basis consisting of integral row-weight vectors.}
\label{claim:WRP4}
There is a finite $O$-basis $(b_{e,j})_{j\in J_e}$ of $L_e$
whose elements have diagonal weights
$\mu_{e,j}\in\mathbb Z^n$ for the action
$\rho_e$.  Under the specified embedding
$K\hookrightarrow\mathbb C$, it is a basis of the literal
$\operatorname{nablaSq}_{\mathbb C}(n,\lambda_e)$.
\begin{proof}
The ambient polynomial module decomposes as the direct sum
of its row-multidegree spaces.  Under the determinant twist,
their torus weights are the row multidegrees plus $s_e$ in
every coordinate.  Only finitely many occur.  Their
projections preserve $\mathcal N_{\nu_e}(K)$: this space
is torus-stable, and over the infinite field $K$ one may
choose a diagonal element on which the distinct occurring
Laurent characters have distinct values and interpolate
the weight projections by polynomials in its action.
Such a diagonal element exists because a nonzero finite
product of Laurent polynomials cannot vanish on the whole
torus of an infinite field.  On the ambient polynomial
module these projections just retain selected monomials,
so they preserve $P_O$ as well.  Consequently
\[
                 L_e=\bigoplus_\mu
                         (L_e\cap\mathcal N_{\nu_e}(K)_\mu).
\]
Each summand is finite free by the same DVR module argument
as WRP3.  Choose a basis in each and take their union.
HP1 identifies extension of the polynomial kernel from
$K$ to $\mathbb C$ coefficientwise for this actual embedding,
so the resulting vectors are a complex basis of the exact
source factor.  No new complex representation is introduced.
\end{proof}
% END WRP4

% BEGIN WRP5
\paragraph{WRP5. The unique lowest coordinate and the exact weight inequalities.}
\label{claim:WRP5}
The basis in WRP4 may be indexed with a distinguished $0_e$
such that
\begin{align}
 \mu_{e,0_e,i}&=(\lambda_e)_{n-1-i},\\
 \sum_i a_i\mu_{e,0_e,i}&\leq\sum_i a_i\mu_{e,j,i}
                         &&(a_0\geq\cdots\geq a_{n-1}),\\
 \sum_i(n-1-i)\mu_{e,0_e,i}
                 &<\sum_i(n-1-i)\mu_{e,j,i} &&(j\ne0_e).
 \label{eq:wrp-weights}
\end{align}
For every upper unitriangular $u$, the coefficient of
$b_{e,0_e}$ in $\rho_e(u)b_{e,j}$ is $\delta_{j,0_e}$.
Every weight has total sum $\sum_i(\lambda_e)_i$.
\begin{proof}
HP1 gives the exact highest polynomial $h_e$, fixed by the
upper unitriangular group and of weight $\lambda_e$.
Let $w_0$ be the reversing permutation matrix and put
$v_e=\rho_e(w_0)h_e$.  It is nonzero, fixed by the lower
unitriangular group, and has weight
$\lambda_e^{\mathrm{rev}}$.

The span of the upper-unitriangular orbit of $v_e$ is the
whole space.  To see this without invoking a new
highest-weight theorem, use the dense open Gaussian cell
$N_+TN_-$ of $\operatorname{GL}_n$.  Lower unipotents fix
$v_e$ and torus elements scale it.  A linear functional
vanishing on its $N_+$-orbit therefore vanishes on its orbit
under this entire open cell.  The matrix coefficient
$g\mapsto\alpha(\rho_e(g)v_e)$ is regular on
$\operatorname{GL}_n$: it is a polynomial in matrix entries
and inverse determinant.  Vanishing on the dense open cell
makes it identically zero.  Thus the span of the $N_+$-orbit
equals the span of the full group orbit, which is the whole
space by the absolute irreducibility proved in HP2.

For an upper elementary matrix $u_{ij}(t)$, $i<j$, and a
vector of weight $\mu$, polynomial expansion in $t$ gives
coefficients of weights
$\mu+r(\epsilon_i-\epsilon_j)$, $r\geq0$.  This follows
by conjugating $u_{ij}(t)$ by a diagonal matrix, or directly
from the row substitution in \eqref{eq:wrp-polynomial}.
The constant coefficient is the original vector.  Every
upper unitriangular matrix is an ordered product of such
elementary matrices.  The preceding spanning assertion
therefore proves that every occurring weight has the form
\[
       \mu=\lambda_e^{\mathrm{rev}}
                 +\sum_{i<j}r_{ij}(\epsilon_i-\epsilon_j),
                         \qquad r_{ij}\geq0.
\]
For the height with coefficients $n-1-i$, every nonzero
positive root has strictly positive height $j-i$.  The
height-minimal part of every expanded orbit vector is
therefore a scalar multiple of $v_e$.  Hence the
$\lambda_e^{\mathrm{rev}}$ weight space is one dimensional.
Its intersection with $L_e$ is a rank-one free direct
summand by WRP4; choose its generator as $b_{e,0_e}$.
Every other chosen basis vector has strictly larger height.

For any decreasing $a_i$, each positive root has
$a_i-a_j\geq0$, giving the weak inequality in
\eqref{eq:wrp-weights}; the strictly decreasing height gives
the strict one.  The same elementary expansions show that
upper unipotents can only increase weights.  They therefore
cannot create a lowest coefficient from another basis
vector, and preserve the lowest coefficient of
$b_{e,0_e}$ itself.  This proves the zeroth-row assertion.
Finally every positive root has coordinate sum zero,
giving the common total weight.  Rank zero has just the
constant polynomial and the same assertions with empty
coordinate sums.
\end{proof}
% END WRP5

% BEGIN WRP6
\paragraph{WRP6. Tensor coordinates and the specified algebraic intertwiner.}
\label{claim:WRP6}
The tensor basis
\[
                 b_J=\bigotimes_{e\in\mathcal E}b_{e,J_e},
                   \qquad J\in\prod_e J_e,
\]
is a basis of the actual $\mathcal V_{\mathbb C}$.  Its
distinguished index is $J_0=(0_e)_e$.  There is an integer
$d\geq0$ and a reindexing by $\{0,\ldots,d\}$ sending
zero to $J_0$.  In these coordinates let $\xi(g)$ be the
matrix of the unchanged representation \eqref{eq:wrp-algrep}.
Then $\xi:G\to\operatorname{GL}_{d+1}(\mathbb C)$ is a
representation, and the coordinate isomorphism
\begin{equation}\label{eq:wrp-intertwiner}
 B:\mathbb C^{d+1}\xrightarrow{\sim}\mathcal V_{\mathbb C},
              \qquad B(\xi(g)x)=\rho(g)B(x),
\end{equation}
proves exactly $\mathrm{IsAlgWt}(\iota,\lambda,\xi)$.
\begin{proof}
Tensor products of finite bases give a basis by the
multilinear expansion of elementary tensors and the
coordinate multilinear functionals.  Each factor has the
nonzero highest polynomial, so each basis index set is
nonempty, and the product index contains $J_0$.  Its finite
cardinality is therefore $d+1$ for some $d$.  Reindexing by
a finite ordinal and then a permutation places $J_0$ at
zero.  The matrix of a group representation in a basis is
invertible, with inverse the matrix at $g^{-1}$; composition
of matrices is exactly composition of operators.
The map $B$ sends standard coordinates to the displayed
basis, so the intertwining identity is its definition.
It uses precisely the maps $\psi_e=\iota\circ j$ and the
source tensor factors.  For empty $\mathcal E$ the tensor
product is $\mathbb C$, the basis has one element, and the
representation is trivial.
\end{proof}
% END WRP6

\subsection{The ten original weight-representation assertions}

% BEGIN WRP7
\paragraph{WRP7. The original lowest-weight character and both zeroth rows.}
\label{claim:WRP7}
For $t\in\mathrm{iwS}.T_f$ write its components uniquely as
$t_v=\operatorname{diag}(y_{v,0},\ldots,y_{v,n-1})$ and define
\begin{equation}\label{eq:wrp-character}
 \chi(t)=\prod_{e=(v,j)\in\mathcal E}\prod_i
                  \psi_e(y_{v,i})^{(\lambda_e)_{n-1-i}}
                  \quad\in\mathbb C^\times.
\end{equation}
This is a group character.  With the coordinates of WRP6,
\begin{align}
 \xi(u)_{0J}&=\delta_{J,J_0} &&(u\in\mathrm{iwS}.N_f),\\
 \xi(t)_{0J}&=\chi(t)\delta_{J,J_0}
                                      &&(t\in\mathrm{iwS}.T_f).
 \label{eq:wrp-rows}
\end{align}
Moreover its value is exactly the character specified by
the source endpoint:
\begin{equation}\label{eq:wrp-lowW}
 \chi(t)=\iota\left(\prod_{v\in V}\prod_i
     \operatorname{algPow}(S(v),\operatorname{lowW}(\lambda,i),
                                                       y_{v,i})\right).
\end{equation}
\begin{proof}
Every diagonal entry $y_{v,i}$ is nonzero, so the products
and their integer powers are nonzero.  Componentwise
multiplication of diagonal matrices proves multiplicativity.
Each tensor factor is diagonal in the weight basis: its
diagonal entry at $j_e$ is
$\prod_i\psi_e(y_{v,i})^{\mu_{e,j_e,i}}$.  Consequently
the whole tensor matrix is diagonal, and its entry at
$J_0$ is \eqref{eq:wrp-character}, by WRP5.  This gives
the torus row.  For an upper unitriangular tuple, the entry
in the $J_0$ row is the product of the corresponding lowest
rows of the factors, hence
$\prod_e\delta_{J_e,0_e}=\delta_{J,J_0}$ by WRP5.

The literal source definition is
\[
 \operatorname{algPow}(w,h,y)=
       \prod_{j:F_w\hookrightarrow\overline{\mathbb Q}_\ell
                         \text{ continuous}} j(y)^{h(j|_L)},
 \qquad
 \operatorname{lowW}(\lambda,i)(\tau)=\lambda(\tau)_{n-1-i}.
\]
The source writes a finite-support product; WRP1 proves that
the entire indexing set is finite, so it is the displayed
finite product.  Expand these definitions, commute the finite
products over $v$, $i$, and $j$, and apply multiplicativity
of $\iota$ and its compatibility with integer powers of
nonzero elements.  The result is exactly
\eqref{eq:wrp-character}.  Thus the character is identified
on every original torus element, with the same local
embeddings and the same field isomorphism.
\end{proof}
% END WRP7

% BEGIN WRP8
\paragraph{WRP8. Integral matrices on the three compact subgroups.}
\label{claim:WRP8}
Every entry of $\xi(g)$ is in $\iota(O)$ if
$g\in\mathrm{iwS}.N$, $g\in\mathrm{iwS}.T$, or
$g\in\mathrm{iwS}.N_m(0)$.  We may therefore take the
original field $c_0$ of $\mathrm{WtRep}$ to be zero.
\begin{proof}
OST26--OST27 identify each component of these three groups
as an element of $\operatorname{GL}_n(\mathcal O_{F_v})$.
The map $\phi_e$ of WRP2 carries both this matrix and its
inverse to integral matrices over $O$, so it gives an
element of $\operatorname{GL}_n(O)$.  WRP3 says that this
element preserves $L_e$.  Its matrix in the $O$-basis
of WRP4 consequently has all entries in $O$.  The image
of this matrix under $\iota$ is the matrix of
$\rho_e(\psi_e(g_v))$ in the complex basis, because all
substitutions and determinant powers commute with the
given field maps.  Entries of the tensor matrix are
products of the entries of these factor matrices, so
they also lie in $\iota(O)$.  This proves, separately,
the fields \texttt{int\_N},
\texttt{integral\_torusCharacter}, and \texttt{int\_Nm}.
The middle field concerns every matrix entry on the compact
torus, as its original definition requires.
\end{proof}
% END WRP8

% BEGIN WRP9
\paragraph{WRP9. Positive normalization and strict contraction outside column zero.}
\label{claim:WRP9}
For every $a\in\mathrm{iwS}.P$, every entry of
$\chi(a)^{-1}\xi(a)$ is in $\iota(O)$.  For the literal
distinguished $z$, every entry of this normalized matrix
outside column zero is in $\iota(\mathfrak m_O)$.
\begin{proof}
Use the actual exponents
$a_v=\operatorname{diag}(\varpi_v^{a_{v,i}})_i$ from
the setup, with each $a_{v,i}$ decreasing in $i$.
The matrix is diagonal by WRP4.  For the diagonal position
$J=(J_e)_e$ put
\[
 \Delta_e(J)=\sum_i a_{v,i}
                \bigl(\mu_{e,J_e,i}-(\lambda_e)_{n-1-i}\bigr)
                \quad(e=(v,j)).
\]
This integer is nonnegative by WRP5.  Using nonzero
$\varpi_v$ and the ordinary rules for integer exponents,
the normalized diagonal entry is the image under $\iota$
of the actual element
\begin{equation}\label{eq:wrp-normalized}
          r_J=\prod_{e=(v,j)}\phi_e(\varpi_v)^{\Delta_e(J)}
                                                        \in O.
\end{equation}
Here $\varpi_v$ is regarded as an element of the local
integer ring, and WRP2 applies.  Nonnegative exponents
are crucial: the normalization has canceled any negative
powers in the original highest weight.  Off-diagonal
entries are zero, proving \texttt{int\_P}.

For $z$ the coefficients are $a_{v,i}=n-1-i$.
If $J\ne J_0$, some $J_e\ne0_e$, so the strict inequality
in WRP5 gives $\Delta_e(J)>0$.  The element
$\phi_e(\varpi_v)$ lies in $\mathfrak m_O$ by WRP2,
and all other factors of \eqref{eq:wrp-normalized} lie
in $O$.  Thus $r_J\in\mathfrak m_O$.  In that column
every off-diagonal entry is zero and also lies in the
maximal ideal.  These elements are the explicit witnesses
required in \texttt{nilpotent\_central\_action}; no
passage to eigenvalues or residue-field coordinates is
being substituted for the required entrywise assertion.
If there is no nonzero column index, the assertion is
vacuous, as in the original definition.
\end{proof}
% END WRP9

% BEGIN WRP10
\paragraph{WRP10. Both integral congruence assertions at every precision.}
\label{claim:WRP10}
For every $k\geq0$ one may choose depth $c=k$ so that
\begin{align}
 \xi(y)_{IJ}&=\delta_{I,J}+\iota(x_{IJ}),
       &x_{IJ}\in\mathfrak m_O^k
                  &&(y\in\mathrm{iwS}.N_m(c)),\\
 \chi(t)&=1+\iota(x),
       &x\in\mathfrak m_O^k
                  &&(t\in\mathrm{iwS}.T_f,\ t\in\mathrm{iwS}.T_c(c)).
 \label{eq:wrp-continuity}
\end{align}
These are exactly \texttt{cont\_Nm} and
\texttt{continuous\_torusCharacter}.
\begin{proof}
By the original definitions checked in OST26--OST27,
each component of either $N_m(c)$ or $T_c(c)$ is
congruent to the identity modulo
$\mathfrak m_{F_v}^{c+1}$, and hence modulo
$\mathfrak m_{F_v}^k$ for $c=k$.  Apply $\phi_e$ and
WRP2 to obtain an integral matrix congruent to the
identity modulo $\mathfrak m_O^k$.  WRP3 gives the same
congruence for its action on $L_e$.  In the factor basis
its entries therefore have the form
$\delta_{i_e,j_e}+u_{i_ej_e}$ with
$u_{i_ej_e}\in\mathfrak m_O^k$.

Entries in the tensor basis are products of these
expressions.  On reducing modulo $\mathfrak m_O^k$,
their product is $\prod_e\delta_{I_e,J_e}=\delta_{I,J}$.
Equivalently their difference from $\delta_{I,J}$ is
an element of that ideal: expand the finite product,
and each other term contains at least one of the
$u_{i_ej_e}$.  Mapping to $\mathbb C$ by $\iota$ gives
the first formula, with witnesses in the original ideal
power rather than in a field topology.  For the torus,
the $00$ entry equals $\chi(t)$ by WRP7; take that entry
of the same congruence.  This gives the second formula.
The argument also covers $k=0$, for which the zeroth
ideal power is $O$, and an empty tensor product has
identity matrix and zero error.  No continuity assertion
about $\iota$ was used.
\end{proof}
% END WRP10

% BEGIN WRP11
\paragraph{WRP11. The exact scalar action on the original center.}
\label{claim:WRP11}
If $t\in\mathrm{iwS}.T_f$ belongs to the center of $G$,
then
\[
                 \xi(t)_{IJ}=\chi(t)\delta_{I,J}.
\]
\begin{proof}
For $n\geq2$, a diagonal matrix commuting with every
elementary matrix has all its diagonal entries equal:
compare the $(i,j)$ entry of its product with
$1+E_{ij}$ in the two orders.  Therefore every component
$t_v$ is a scalar matrix, since it commutes with the
elementary matrices supported in that component of $G$.
For $n=1$ it is already scalar; for $n=0$ all groups
and representation factors are trivial.  Write the
scalar as $s_v$ when $n>0$.

By WRP5 every weight of factor $e=(v,j)$ has the same
sum $\sum_i(\lambda_e)_i$.  Its diagonal eigenvalue
at $s_v$ is consequently
$\psi_e(s_v)^{\sum_i(\lambda_e)_i}$, independent of
the basis index.  Each factor is thus scalar, and
their product equals its own lowest-coordinate value,
which is $\chi(t)$ by WRP7.  This proves the entrywise
\texttt{central} field, including all off-diagonal
zeros, for the actual full product center.
\end{proof}
% END WRP11

% BEGIN WRP12
\paragraph{WRP12. Installation of the exact original \texttt{wtRep} endpoint.}
\label{claim:WRP12}
For every original $(L,\ell,n,V,S,\iota,\lambda,c)$,
WRP1--WRP11 construct $c'\geq c$ and a representation
\[
                  R:\mathrm{WtRep}(\mathrm{iwS}(n,S),c'.O,\mathbb C)
\]
such that $\mathrm{IsAlgWt}(\iota,\lambda,R.\xi)$ and
the equality \eqref{eq:wrp-lowW} hold for every original
torus tuple.  This is the exact statement of
\texttt{UpperTheory.V.Hi.wtRep}, in
\texttt{ProdZpowConst.lean}, at the pinned source revision.
\begin{proof}
Use the subalgebra $c'.O$ constructed in WRP2, with
complex algebra map equal to the restriction of the
specified $\iota$.  Take $d,\xi$ from WRP6,
$\chi$ from WRP7, and $c_0=0$ from WRP8.  The ten
original assertion fields are supplied individually
by the following proved statements:
\begin{center}
\begin{tabular}{ll}
Original field & Written proof \\
\hline
\texttt{int\_N} & WRP8 \\
\texttt{integral\_torusCharacter} & WRP8 \\
\texttt{int\_Nm} & WRP8, with $c_0=0$ \\
\texttt{int\_P} & WRP9 \\
\texttt{row\_N} & WRP7 \\
\texttt{torusAction} & WRP7 \\
\texttt{nilpotent\_central\_action} & WRP9 \\
\texttt{cont\_Nm} & WRP10, every $k$ \\
\texttt{continuous\_torusCharacter} & WRP10, every $k$ \\
\texttt{central} & WRP11
\end{tabular}
\end{center}
WRP6 supplies the explicit linear intertwiner $B$ with
the literal original tensor representation, and WRP7
proves the required exact character identity.  These
are precisely the two additional conclusions of
the source existential theorem.  The construction
allows every prime $\ell$, every dominant integer
weight, every given algebraic isomorphism $\iota$,
and every original coefficient datum $c$.  It retains
the cases $n=0$, empty $V$, and places having no
continuous embedding into $\overline{\mathbb Q}_\ell$.
\end{proof}
% END WRP12

\subsection{Scope of the completed supplier}

The basis used here need not be the particular basis
selected by the source's auxiliary \texttt{sqBasis}
theorem over arbitrary domains.  The endpoint just
proved is existential, and its identification with
\texttt{algRep} is exactly the explicit intertwiner
\eqref{eq:wrp-intertwiner}; every source assertion holds
for this witness with the unchanged character
\eqref{eq:wrp-lowW}.  Thus this permitted simplification
does not change the representation, the coefficient
embedding, the local groups, the normalized positive
operators, or the input weight.

In conjunction with OCT1--OCT19 this supplies the
previously explicit \texttt{WtRep} hypothesis of the
ordinary control theorem for each of the actual
algebraic weights and the fixed complex embedding.
The coefficient extension is the constructed inclusion
$c.O\subset c'.O$, so subsequent use must still make
the corresponding original coefficient map explicit.
This chapter by itself does not assert the entire
\texttt{SpecIntIff.spec} theorem: the finite-order torus
character chosen there, the single-map assembly over
all levels, and the arithmetic specialization contract
must be supplied where they are used.  Nor does it
assert a completed same-ideal Hida construction or
the global RH/GRH obligation.  Its completed output
is the exact quantified \texttt{wtRep} supplier above.
